\documentclass[a4paper,10pt]{article}

\usepackage{geometry}
\usepackage{fontspec}
\usepackage{polyglossia}
\setmainlanguage{russian}
\setotherlanguage{english}
\usepackage{cv-simple}
% Russian CVs: TeX Gyre Termes (Times-like, Cyrillic + Latin) to match the English design
\newcommand{\cvfontsru}{%
  \setmainfont{Times New Roman}[
    BoldFont = {Times New Roman Bold},
    ItalicFont = {Times New Roman Italic},
    BoldItalicFont = {Times New Roman Bold Italic},
    Ligatures = TeX,
  ]%
  \newfontfamily\cyrillicfont{Times New Roman}[
    BoldFont = {Times New Roman Bold},
    ItalicFont = {Times New Roman Italic},
    BoldItalicFont = {Times New Roman Bold Italic},
    Ligatures = TeX,
  ]%
  \newfontfamily\englishfont{Times New Roman}[
    BoldFont = {Times New Roman Bold},
    ItalicFont = {Times New Roman Italic},
    BoldItalicFont = {Times New Roman Bold Italic},
    Ligatures = TeX,
  ]%
}

\newcommand{\cvpagemargins}{%
  \geometry{a4paper,margin=0.55in,top=0.15in,bottom=0.08in}%
  \renewcommand{\baselinestretch}{0.88}%
}

\cvfontsru
\cvpagemargins

\begin{document}

\cvname{Наталья А. Агапова}{ML Engineer}
\cvcontactsru

\section{Образование}
\resumeSubHeadingListStart
  \resumeSubheading{Иннополис, Россия}{2022 – 2026}{Университет Иннополис}{}
  \resumeItemListStart
    \resumeItem{\textbf{Специальность:} Бакалавриат «Прикладной искусственный интеллект» (средний балл: 4.02 из 5)}
  \resumeItemListEnd
\resumeSubHeadingListEnd

\section{Опыт}
\resumeSubHeadingListStart

  \resumeProject{\href{https://croissanstudio.ru}{Croissan Studio}}{\textbf{апр. 2025 – н.в.}}
  \resumeRoleListStart
    \resumeRole{ML Software Engineer}{апр. 2026 – н.в.}
    \resumeItemListStart
      \resumeItem{Работаю ML-инженером в клиентских командах, которые делают AI-продукты от начала до конца, вместе с дизайнерами и инженерами.}
      \resumeItem{Полностью дообучала языковые модели и адаптировала другие через промптинг для клиентских продуктов.}
      \resumeItem{Сравнивала версии моделей по метрикам задачи, чтобы выбрать, что идёт в продукт.}
      \resumeItem{Обучала и адаптировала генеративные модели изображений и говорящих аватаров для клиентских проектов.}
    \resumeItemListEnd

    \resumeRole{ML Engineer \& Product Designer}{июль 2025 – апр. 2026}
    \resumeItemListStart
      \resumeItem{Переработала базовую модель NeuroArtist, свободно выбирая и тестируя собственный подход.}
      \resumeItem{Сравнивала результаты модели с требованиями клиентов к качеству, выбирая, что идёт в релиз.}
      \resumeItem{Также руководила ребрендингом NeuroArtist (физические и цифровые материалы) и участвовала в редизайне сайта Croissan Studio.}
    \resumeItemListEnd

    \resumeRole{Software Engineer Intern}{апр. – июль 2025}
    \resumeItemListStart
      \resumeItem{Работала над сайтом NeuroArtist.}
      \resumeItem{Помогала с ML-задачами в команде, после стажировки получила предложение на роль ML Engineer.}
      \resumeItem{Участвовала в клиентских проектах, например туристический сайт и мероприятие «Москва 2030» и др.}
    \resumeItemListEnd
  \resumeRoleListEnd
  \resumeItemListStart
    \cvtech{Применила знания}{ в Python, PyTorch, Hugging Face Transformers, Diffusers и LoRA fine-tuning.}
    \cvletter{https://drive.google.com/file/d/1VtzJ118F98tqOT2DnjbEX7sI2dEZRsJv/view?usp=sharing}{Рекомендательное письмо от CEO студии}
  \resumeItemListEnd

  \resumeProject{\href{https://asimovlab.ru}{Asimov Lab} (startup)}{\textbf{фев. – сент. 2024}}
  \resumeRoleListStart
    \resumeRole{Frontend Team Lead \& ML Developer}{}
    \resumeItemListStart
      \resumeItem{Привела фронтенд-команду к MVP AI-платформы, которая генерирует тесты для преподавателей и студентов.}
      \resumeItem{Собрала MVP-версию сайта от начала до конца.}
      \resumeItem{Один из шести ML-разработчиков, стоявших за ключевой моделью продукта.}
    \resumeItemListEnd
  \resumeRoleListEnd
  \resumeItemListStart
    \cvtech{Применила знания}{ в React, TypeScript, Python и scikit-learn.}
  \resumeItemListEnd

  \resumeProject{Фриланс}{\textbf{2022 – н.в.}}
  \resumeRoleListStart
    \resumeRole{ML Engineer}{}
    \resumeItemListStart
      \resumeItem{Прикладные ML-задачи для независимых клиентов.}
    \resumeItemListEnd
  \resumeRoleListEnd
  \resumeItemListStart
    \cvtech{Применила знания}{ в Python, PyTorch и scikit-learn.}
  \resumeItemListEnd

\resumeSubHeadingListEnd

\section{Личные проекты}
\resumeSubHeadingListStart

  \resumeProject{\href{https://github.com/natagapova/resume-screening}{Fair Resume Screening}}{NLP · Explainability · Debiasing · Публикация, первый автор}
  \resumeItemListStart
    \resumeItem{Обучила BERT сортировать резюме по 9 IT-категориям на реальных данных HeadHunter, затем проверила модель на предвзятость и протестировала 39+ способов её снизить.}
    \resumeItem{Выяснила, что замена только города кандидата меняла решение модели в 7,7\% случаев (baseline accuracy 60,9\%) -- опубликовала эту находку как первый автор.}
    \cvtech{Использовано:}{ Python, PyTorch, BERT, Captum, pandas.}
  \resumeItemListEnd

  \resumeProject{\href{https://github.com/natagapova/emotion-detection}{Emotion Detection}}{Computer Vision · On-device Deployment}
  \resumeItemListStart
    \resumeItem{Обучила CNN распознавать эмоции по лицу с accuracy 58,7\% (до 60,8\% с ансамблированием), затем сжала модель для работы на устройстве с задержкой всего 1,15\,мс (Core ML, ONNX Runtime Web).}
    \cvtech{Использовано:}{ Python, PyTorch, Core ML, ONNX Runtime.}
  \resumeItemListEnd

  \resumeProject{\href{https://github.com/natagapova/personal-knowledge-system}{Personal Knowledge System}}{RAG · NLP · Embeddings}
  \resumeItemListStart
    \resumeItem{Собрала личного AI-ассистента, который отвечает только по вашим PDF и всегда указывает файл и страницу -- на локальной LLM (Llama 3.2), так что ничего не уходит с вашего устройства.}
    \cvtech{Использовано:}{ Python, SentenceTransformer, ChromaDB, Ollama.}
  \resumeItemListEnd

\resumeSubHeadingListEnd

\section{Навыки}
\cvskillsblock{%
\textbf{ML / AI:} \textit{(свободно)}: Python, PyTorch, BERT, Captum, оценка и debiasing моделей, RAG (ChromaDB, SentenceTransformer) \quad \textit{(знакома)}: LLM (Ollama), OpenCV, MediaPipe, Core ML, ONNX Runtime\\
\textbf{Инструменты:} \textit{(знакома)}: Figma, фронтенд-разработка, Git/GitHub\\
\cvlanguagesru}

\section{Публикации}
\cvpublication{Agapova, N. A.; Lukmanov, R. A. ``Improving Fairness in AI-Powered Recruitment: An Interpretable Resume Screening System'', 2026}{}{10.66693/mathai.1017}
\cvpublication{Kholodov, Y. A.; Salloum, H.; Agapova, N. A. ``Quantum Annealing Approaches to Breaking RSA Encryption'', \textit{Voprosy kiberbezopasnosti}, 2025}{}{10.21681/2311-3456-2025-4-127-133}

\end{document}
